% Generated draft: requires mathematical and source review.
\documentclass[11pt]{article}
\usepackage[margin=25mm]{geometry}
\usepackage{fontspec}
\setmainfont{Latin Modern Roman}
\usepackage{amsmath,amssymb,mathtools,unicode-math}
\setmathfont{Latin Modern Math}
\usepackage{enumitem,needspace,xurl,hyperref}
\setlength{\parindent}{0pt}
\setlength{\parskip}{5pt}
\setlength{\emergencystretch}{3em}
\newcommand{\N}{\mathbb{N}}
\newcommand{\Z}{\mathbb{Z}}
\newcommand{\Q}{\mathbb{Q}}
\newcommand{\R}{\mathbb{R}}
\newcommand{\C}{\mathbb{C}}
\newcommand{\F}{\mathbb{F}}
\newcommand{\A}{\mathbb{A}}
\newcommand{\PP}{\mathbb{P}}
\newcommand{\E}{\mathbb{E}}
\newcommand{\eps}{\varepsilon}
\newcommand{\dd}{\,\mathrm{d}}
\begin{document}
\section*{UnsolvedMath ID: 700453}
\textbf{Source problem code:} catalogue entry
\subsection{Definitions and mathematical statement}
Let \(F\) be a field and \(G\) a group, and let \(FG\) denote their group algebra. An irreducible \(FG\)-module means a nonzero unital left \(FG\)-module with no proper nonzero submodules. For each such module \(M\), put
\[
D_M=\operatorname{End}_{FG}(M),\qquad
K_M=Z(D_M)=\{c\in D_M:cd=dc\text{ for every }d\in D_M\}.
\]
By Schur's lemma, \(D_M\) is a division ring and \(K_M\) is a field. Regard \(M\) as a vector space over \(K_M\) through the action \(c\cdot m=c(m)\).

Characterize all pairs \((G,F)\) satisfying
\[
\text{for every irreducible }FG\text{-module }M,
\qquad \dim_{K_M}M<\infty.
\]
The requested characterization allows arbitrary groups and fields. As a more restricted target, characterize all such pairs with \(G\) soluble: writing \(G^{(0)}=G\) and \(G^{(n+1)}=[G^{(n)},G^{(n)}]\), this means that \(G^{(r)}=\{1\}\) for some integer \(r\geq 0\). No common finite bound on the dimensions is required. The dimensions are taken over \(K_M\), not merely over \(F\).

The supplied source refers to partial answers in the Wehrfritz paper listed below, but does not specify their hypotheses or conclusions; no further restriction of the requested characterization can be extracted from that reference alone.
\subsection{Short English statement}
Determine exactly which groups \(G\) and fields \(F\) have the property that every irreducible representation over \(F\) is finite dimensional over the centre of its ring of module endomorphisms. The source also requests, as a restricted target, a characterization for soluble groups.
\subsection{Sources}
\begin{itemize}
\item Local UnsolvedMath registry, supplied complete source record: permanent ID \(700453\), exact source code \texttt{catalogue entry}, title \emph{Kourovka Notebook Problem 8.1}; upstream revision \texttt{d371ee917ac9647fbb06cd6fc529deec6f3dda82}. The supplied registry status is \texttt{open}; this status has not been independently verified.
\item \emph{The Kourovka Notebook (Unsolved Problems in Group Theory)}, original problem 8.1, 8th issue (1982), proposed by B. A. F. Wehrfritz. The supplied background identifies the consulted edition as the 21st issue edition (2026), \texttt{arXiv:1401.0300v48}. The collection description names E. I. Khukhro and V. D. Mazurov as editors.
\item Reference for partial answers, reproduced as supplied: B. A. F. Wehrfritz, \emph{Glasgow Math. J.}, \textbf{24}, no. 1 (1983), 169--176. The supplied record gives no article title or link and does not state the partial results.
\end{itemize}
\end{document}
